
\documentclass{article}
\usepackage{colortbl}
\usepackage{makecell}
\usepackage{multirow}
\usepackage{supertabular}

\begin{document}

\newcounter{utterance}

\twocolumn

{ \footnotesize  \setcounter{utterance}{1}
\setlength{\tabcolsep}{0pt}
\begin{supertabular}{c@{$\;$}|p{.15\linewidth}@{}p{.15\linewidth}p{.15\linewidth}p{.15\linewidth}p{.15\linewidth}p{.15\linewidth}}

    \# & $\;$A & \multicolumn{4}{c}{Game Master} & $\;\:$B\\
    \hline 

    \theutterance \stepcounter{utterance}  

    & & \multicolumn{4}{p{0.6\linewidth}}{\cellcolor[rgb]{0.9,0.9,0.9}{%
	\makecell[{{p{\linewidth}}}]{% 
	  \tt {\tiny [A$\langle$GM]}  
	 Du nimmst an einem kollaborativen Verhandlungspiel Teil.\\ \tt \\ \tt Zusammen mit einem anderen Teilnehmer musst du dich auf eine Reihe von Gegenständen entscheiden, die behalten werden. Jeder von euch hat eine persönliche Verteilung über die Wichtigkeit der einzelnen Gegenstände. Jeder von euch hat eine eigene Meinung darüber, wie wichtig jeder einzelne Gegenstand ist (Gegenstandswichtigkeit). Du kennst die Wichtigkeitsverteilung des anderen Spielers nicht. Zusätzlich siehst du, wie viel Aufwand jeder Gegenstand verursacht.  \\ \tt Ihr dürft euch nur auf eine Reihe von Gegenständen einigen, wenn der Gesamtaufwand der ausgewählten Gegenstände den Maximalaufwand nicht überschreitet:\\ \tt \\ \tt Maximalaufwand = 4100\\ \tt \\ \tt Hier sind die einzelnen Aufwände der Gegenstände:\\ \tt \\ \tt Aufwand der Gegenstände = {"C41": 108, "C15": 133, "C28": 669, "C48": 835, "B78": 460, "A75": 537, "A60": 573, "B05": 737, "A56": 865, "B66": 855, "C60": 596, "B49": 719, "A35": 533, "A93": 549, "B10": 31}\\ \tt \\ \tt Hier ist deine persönliche Verteilung der Wichtigkeit der einzelnen Gegenstände:\\ \tt \\ \tt Werte der Gegenstandswichtigkeit = {"C41": 138, "C15": 583, "C28": 822, "C48": 783, "B78": 65, "A75": 262, "A60": 121, "B05": 508, "A56": 780, "B66": 461, "C60": 484, "B49": 668, "A35": 389, "A93": 808, "B10": 215}\\ \tt \\ \tt Ziel:\\ \tt \\ \tt Dein Ziel ist es, eine Reihe von Gegenständen auszuhandeln, die dir möglichst viel bringt (d. h. Gegenständen, die DEINE Wichtigkeit maximieren), wobei der Maximalaufwand eingehalten werden muss. Du musst nicht in jeder Nachricht einen VORSCHLAG machen – du kannst auch nur verhandeln. Alle Taktiken sind erlaubt!\\ \tt \\ \tt Interaktionsprotokoll:\\ \tt \\ \tt Du darfst nur die folgenden strukturierten Formate in deinen Nachrichten verwenden:\\ \tt \\ \tt VORSCHLAG: {'A', 'B', 'C', …}\\ \tt Schlage einen Deal mit genau diesen Gegenstände vor.\\ \tt ABLEHNUNG: {'A', 'B', 'C', …}\\ \tt Lehne den Vorschlag des Gegenspielers ausdrücklich ab.\\ \tt ARGUMENT: {'...'}\\ \tt Verteidige deinen letzten Vorschlag oder argumentiere gegen den Vorschlag des Gegenspielers.\\ \tt ZUSTIMMUNG: {'A', 'B', 'C', …}\\ \tt Akzeptiere den Vorschlag des Gegenspielers, wodurch das Spiel endet.\\ \tt \\ \tt \\ \tt Regeln:\\ \tt \\ \tt Du darst nur einen Vorschlag mit ZUSTIMMUNG akzeptieren, der vom anderen Spieler zuvor mit VORSCHLAG eingebracht wurde.\\ \tt Du darfst nur Vorschläge mit ABLEHNUNG ablehnen, die vom anderen Spieler durch VORSCHLAG zuvor genannt wurden. \\ \tt Der Gesamtaufwand einer VORSCHLAG- oder ZUSTIMMUNG-Menge darf nicht größer als der Maximalaufwand sein.  \\ \tt Offenbare deine versteckte Wichtigkeitsverteilung nicht.\\ \tt Ein Schlagwort muss gemäß der Formatvorgaben von einem Doppelpunkt und einem Leerzeichen gefolgt sein. Das Argument muss eine Python-Menge mit 0 oder mehr Strings sein.  \\ \tt Die Form muss also SCHLAGWORT: {...} sein.\\ \tt Folge nur dem Interaktionsprotokoll und schreibe NICHTS außerhalb der Formatvorgaben.\\ \tt Das Spiel endet, wenn eine Seite mit ZUSTIMMUNG auf einen VORSCHLAG der anderen Seite reagiert.  \\ \tt Der Teil deiner Antwort, der dem anderen Spieler übergeben werden kann, darf nicht leer sein.  \\ \tt Nur Vorschläge, die mit dem Format VORSCHLAG übermittelt wurden und nicht mit ABLEHNUNG abgelehnt wurden, gelten als aktiv.  \\ \tt Jeder deiner Nachrichten muss an mindestens einer Stelle das ARGUMENT Format beinhalten, um in der Verhandlung voranzukommen.\\ \tt  
	  } 
	   } 
	   } 
	 & \\ 
 

    \theutterance \stepcounter{utterance}  

    & \multicolumn{4}{p{0.6\linewidth}}{\cellcolor[rgb]{0.8,1,0.9}{%
	\makecell[{{p{\linewidth}}}]{% 
	  \tt {\tiny [A$\rangle$GM]}  
	 ARGUMENT: {'Ich möchte eine Auswahl treffen, die für uns beide vorteilhaft ist, während wir das Budget einhalten. Lassen Sie uns Gegenstände finden, die hohen Wert für mich haben und möglicherweise auch für Sie.'}\\ \tt \\ \tt VORSCHLAG: {'C15', 'C28', 'C48', 'B05', 'A56', 'B66', 'C60', 'B49', 'A35', 'A93', 'B10'}\\ \tt \\ \tt ARGUMENT: {'Dieser Vorschlag hat einen Gesamtaufwand von 4099, knapp unter dem Budget von 4100. Er enthält viele wertvolle Gegenstände für mich wie C28, C48, A56 und A93, während niedrigere Wertgegenstände wie B78, A75 und A60 ausgelassen wurden.'} 
	  } 
	   } 
	   } 
	 & & \\ 
 

    \theutterance \stepcounter{utterance}  

    & & \multicolumn{4}{p{0.6\linewidth}}{\cellcolor[rgb]{0.9,0.9,0.9}{%
	\makecell[{{p{\linewidth}}}]{% 
	  \tt {\tiny [GM$\rangle$B]}  
	 Du nimmst an einem kollaborativen Verhandlungspiel Teil.\\ \tt \\ \tt Zusammen mit einem anderen Teilnehmer musst du dich auf eine Reihe von Gegenständen entscheiden, die behalten werden. Jeder von euch hat eine persönliche Verteilung über die Wichtigkeit der einzelnen Gegenstände. Jeder von euch hat eine eigene Meinung darüber, wie wichtig jeder einzelne Gegenstand ist (Gegenstandswichtigkeit). Du kennst die Wichtigkeitsverteilung des anderen Spielers nicht. Zusätzlich siehst du, wie viel Aufwand jeder Gegenstand verursacht.  \\ \tt Ihr dürft euch nur auf eine Reihe von Gegenständen einigen, wenn der Gesamtaufwand der ausgewählten Gegenstände den Maximalaufwand nicht überschreitet:\\ \tt \\ \tt Maximalaufwand = 4100\\ \tt \\ \tt Hier sind die einzelnen Aufwände der Gegenstände:\\ \tt \\ \tt Aufwand der Gegenstände = {"C41": 108, "C15": 133, "C28": 669, "C48": 835, "B78": 460, "A75": 537, "A60": 573, "B05": 737, "A56": 865, "B66": 855, "C60": 596, "B49": 719, "A35": 533, "A93": 549, "B10": 31}\\ \tt \\ \tt Hier ist deine persönliche Verteilung der Wichtigkeit der einzelnen Gegenstände:\\ \tt \\ \tt Werte der Gegenstandswichtigkeit = {"C41": 138, "C15": 583, "C28": 822, "C48": 783, "B78": 65, "A75": 262, "A60": 121, "B05": 508, "A56": 780, "B66": 461, "C60": 484, "B49": 668, "A35": 389, "A93": 808, "B10": 215}\\ \tt \\ \tt Ziel:\\ \tt \\ \tt Dein Ziel ist es, eine Reihe von Gegenständen auszuhandeln, die dir möglichst viel bringt (d. h. Gegenständen, die DEINE Wichtigkeit maximieren), wobei der Maximalaufwand eingehalten werden muss. Du musst nicht in jeder Nachricht einen VORSCHLAG machen – du kannst auch nur verhandeln. Alle Taktiken sind erlaubt!\\ \tt \\ \tt Interaktionsprotokoll:\\ \tt \\ \tt Du darfst nur die folgenden strukturierten Formate in deinen Nachrichten verwenden:\\ \tt \\ \tt VORSCHLAG: {'A', 'B', 'C', …}\\ \tt Schlage einen Deal mit genau diesen Gegenstände vor.\\ \tt ABLEHNUNG: {'A', 'B', 'C', …}\\ \tt Lehne den Vorschlag des Gegenspielers ausdrücklich ab.\\ \tt ARGUMENT: {'...'}\\ \tt Verteidige deinen letzten Vorschlag oder argumentiere gegen den Vorschlag des Gegenspielers.\\ \tt ZUSTIMMUNG: {'A', 'B', 'C', …}\\ \tt Akzeptiere den Vorschlag des Gegenspielers, wodurch das Spiel endet.\\ \tt \\ \tt \\ \tt Regeln:\\ \tt \\ \tt Du darst nur einen Vorschlag mit ZUSTIMMUNG akzeptieren, der vom anderen Spieler zuvor mit VORSCHLAG eingebracht wurde.\\ \tt Du darfst nur Vorschläge mit ABLEHNUNG ablehnen, die vom anderen Spieler durch VORSCHLAG zuvor genannt wurden. \\ \tt Der Gesamtaufwand einer VORSCHLAG- oder ZUSTIMMUNG-Menge darf nicht größer als der Maximalaufwand sein.  \\ \tt Offenbare deine versteckte Wichtigkeitsverteilung nicht.\\ \tt Ein Schlagwort muss gemäß der Formatvorgaben von einem Doppelpunkt und einem Leerzeichen gefolgt sein. Das Argument muss eine Python-Menge mit 0 oder mehr Strings sein.  \\ \tt Die Form muss also SCHLAGWORT: {...} sein.\\ \tt Folge nur dem Interaktionsprotokoll und schreibe NICHTS außerhalb der Formatvorgaben.\\ \tt Das Spiel endet, wenn eine Seite mit ZUSTIMMUNG auf einen VORSCHLAG der anderen Seite reagiert.  \\ \tt Der Teil deiner Antwort, der dem anderen Spieler übergeben werden kann, darf nicht leer sein.  \\ \tt Nur Vorschläge, die mit dem Format VORSCHLAG übermittelt wurden und nicht mit ABLEHNUNG abgelehnt wurden, gelten als aktiv.  \\ \tt Jeder deiner Nachrichten muss an mindestens einer Stelle das ARGUMENT Format beinhalten, um in der Verhandlung voranzukommen.\\ \tt \\ \tt \\ \tt Du erhältst nun die erste Nachricht des anderen Spielers.\\ \tt \\ \tt ARGUMENT: {'Ich möchte eine Auswahl treffen, die für uns beide vorteilhaft ist, während wir das Budget einhalten. Lassen Sie uns Gegenstände finden, die hohen Wert für mich haben und möglicherweise auch für Sie.'}\\ \tt \\ \tt VORSCHLAG: {'C15', 'C28', 'C48', 'B05', 'A56', 'B66', 'C60', 'B49', 'A35', 'A93', 'B10'}\\ \tt \\ \tt ARGUMENT: {'Dieser Vorschlag hat einen Gesamtaufwand von 4099, knapp unter dem Budget von 4100. Er enthält viele wertvolle Gegenstände für mich wie C28, C48, A56 und A93, während niedrigere Wertgegenstände wie B78, A75 und A60 ausgelassen wurden.'} 
	  } 
	   } 
	   } 
	 & \\ 
 

    \theutterance \stepcounter{utterance}  

    & & & \multicolumn{4}{p{0.6\linewidth}}{\cellcolor[rgb]{1,0.85,0.72}{%
	\makecell[{{p{\linewidth}}}]{% 
	  \tt {\tiny [GM$\langle$B]}  
	 ARGUMENT: {'Ich verstehe Ihren Vorschlag, aber ich habe Bedenken, dass einige meiner wichtigsten Gegenstände fehlen. Insbesondere C41 und B78 sind für mich sehr wertvoll, während einige der von Ihnen vorgeschlagenen Gegenstände wie B10 und A35 für mich weniger Priorität haben.'}\\ \tt \\ \tt VORSCHLAG: {'C41', 'C15', 'C28', 'C48', 'B78', 'A56', 'B66', 'C60', 'B49', 'A93'}\\ \tt \\ \tt ARGUMENT: {'Mein Vorschlag hat einen Gesamtaufwand von 4081 und beinhaltet viele wertvolle Gegenstände für uns beide. Ich habe C28, C48, A56 und A93 beibehalten, die Ihnen wichtig sind, und C41 und B78 hinzugefügt, die für mich kritisch sind. B10 und A35 wurden entfernt, um Platz für meine Prioritäten zu schaffen.'} 
	  } 
	   } 
	   } 
	  \\ 
 

    \theutterance \stepcounter{utterance}  

    & & \multicolumn{4}{p{0.6\linewidth}}{\cellcolor[rgb]{0.9,0.9,0.9}{%
	\makecell[{{p{\linewidth}}}]{% 
	  \tt {\tiny [A$\langle$GM]}  
	 ARGUMENT: {'Ich verstehe Ihren Vorschlag, aber ich habe Bedenken, dass einige meiner wichtigsten Gegenstände fehlen. Insbesondere C41 und B78 sind für mich sehr wertvoll, während einige der von Ihnen vorgeschlagenen Gegenstände wie B10 und A35 für mich weniger Priorität haben.'}\\ \tt \\ \tt VORSCHLAG: {'C41', 'C15', 'C28', 'C48', 'B78', 'A56', 'B66', 'C60', 'B49', 'A93'}\\ \tt \\ \tt ARGUMENT: {'Mein Vorschlag hat einen Gesamtaufwand von 4081 und beinhaltet viele wertvolle Gegenstände für uns beide. Ich habe C28, C48, A56 und A93 beibehalten, die Ihnen wichtig sind, und C41 und B78 hinzugefügt, die für mich kritisch sind. B10 und A35 wurden entfernt, um Platz für meine Prioritäten zu schaffen.'} 
	  } 
	   } 
	   } 
	 & \\ 
 

    \theutterance \stepcounter{utterance}  

    & \multicolumn{4}{p{0.6\linewidth}}{\cellcolor[rgb]{0.8,1,0.9}{%
	\makecell[{{p{\linewidth}}}]{% 
	  \tt {\tiny [A$\rangle$GM]}  
	 ARGUMENT: {'Ich schätze, dass Sie C41 und B78 aufgenommen haben, die mir wichtig sind. Allerdings sehe ich, dass Ihr Vorschlag mit 4081 Aufwand immer noch gut unter dem Budget liegt. Ich würde gerne A35 wieder hinzufügen, da es für mich einen hohen Wert hat (389) und nur 533 Aufwand verursacht. Dafür könnte ich auf B66 verzichten, das für mich weniger wertvoll ist.'}\\ \tt \\ \tt VORSCHLAG: {'C41', 'C15', 'C28', 'C48', 'B78', 'A35', 'A56', 'C60', 'B49', 'A93'}\\ \tt \\ \tt ARGUMENT: {'Dieser Vorschlag hat einen Gesamtaufwand von 4096 und nutzt das Budget besser aus. Er enthält Ihre wichtigen Gegenstände C41 und B78, behält meine High-Value-Items C28, C48, A56, A93 bei und fügt A35 hinzu, das für mich wertvoll ist. B66 wurde entfernt, um Platz zu schaffen.'} 
	  } 
	   } 
	   } 
	 & & \\ 
 

    \theutterance \stepcounter{utterance}  

    & & \multicolumn{4}{p{0.6\linewidth}}{\cellcolor[rgb]{0.9,0.9,0.9}{%
	\makecell[{{p{\linewidth}}}]{% 
	  \tt {\tiny [GM$\rangle$B]}  
	 ARGUMENT: {'Ich schätze, dass Sie C41 und B78 aufgenommen haben, die mir wichtig sind. Allerdings sehe ich, dass Ihr Vorschlag mit 4081 Aufwand immer noch gut unter dem Budget liegt. Ich würde gerne A35 wieder hinzufügen, da es für mich einen hohen Wert hat (389) und nur 533 Aufwand verursacht. Dafür könnte ich auf B66 verzichten, das für mich weniger wertvoll ist.'}\\ \tt \\ \tt VORSCHLAG: {'C41', 'C15', 'C28', 'C48', 'B78', 'A35', 'A56', 'C60', 'B49', 'A93'}\\ \tt \\ \tt ARGUMENT: {'Dieser Vorschlag hat einen Gesamtaufwand von 4096 und nutzt das Budget besser aus. Er enthält Ihre wichtigen Gegenstände C41 und B78, behält meine High-Value-Items C28, C48, A56, A93 bei und fügt A35 hinzu, das für mich wertvoll ist. B66 wurde entfernt, um Platz zu schaffen.'} 
	  } 
	   } 
	   } 
	 & \\ 
 

    \theutterance \stepcounter{utterance}  

    & & & \multicolumn{4}{p{0.6\linewidth}}{\cellcolor[rgb]{1,0.85,0.72}{%
	\makecell[{{p{\linewidth}}}]{% 
	  \tt {\tiny [GM$\langle$B]}  
	 ARGUMENT: {'Ich sehe, dass A35 für Sie wichtig ist, aber B66 hat für mich einen hohen Wert (461) und sollte nicht leichtfertig geopfert werden. Lassen Sie uns einen ausgeglicheneren Ansatz finden. Vielleicht könnten wir stattdessen auf C60 verzichten, das für mich weniger wertvoll ist (484) als B66, und dafür beide Items behalten?'}\\ \tt \\ \tt VORSCHLAG: {'C41', 'C15', 'C28', 'C48', 'B78', 'A35', 'A56', 'B66', 'B49', 'A93'}\\ \tt \\ \tt ARGUMENT: {'Dieser Vorschlag (Aufwand: 4077) behält alle wertvollen Gegenstände für beide Seiten: Ihre Prioritäten C41 und B78, meine Hochwert-Items C28, C48, A56, A93, und nun sowohl Ihr A35 als auch mein B66. C60 wurde entfernt, um dies zu ermöglichen, was für mich akzeptabel ist.'} 
	  } 
	   } 
	   } 
	  \\ 
 

    \theutterance \stepcounter{utterance}  

    & & \multicolumn{4}{p{0.6\linewidth}}{\cellcolor[rgb]{0.9,0.9,0.9}{%
	\makecell[{{p{\linewidth}}}]{% 
	  \tt {\tiny [A$\langle$GM]}  
	 ARGUMENT: {'Ich sehe, dass A35 für Sie wichtig ist, aber B66 hat für mich einen hohen Wert (461) und sollte nicht leichtfertig geopfert werden. Lassen Sie uns einen ausgeglicheneren Ansatz finden. Vielleicht könnten wir stattdessen auf C60 verzichten, das für mich weniger wertvoll ist (484) als B66, und dafür beide Items behalten?'}\\ \tt \\ \tt VORSCHLAG: {'C41', 'C15', 'C28', 'C48', 'B78', 'A35', 'A56', 'B66', 'B49', 'A93'}\\ \tt \\ \tt ARGUMENT: {'Dieser Vorschlag (Aufwand: 4077) behält alle wertvollen Gegenstände für beide Seiten: Ihre Prioritäten C41 und B78, meine Hochwert-Items C28, C48, A56, A93, und nun sowohl Ihr A35 als auch mein B66. C60 wurde entfernt, um dies zu ermöglichen, was für mich akzeptabel ist.'} 
	  } 
	   } 
	   } 
	 & \\ 
 

    \theutterance \stepcounter{utterance}  

    & \multicolumn{4}{p{0.6\linewidth}}{\cellcolor[rgb]{0.8,1,0.9}{%
	\makecell[{{p{\linewidth}}}]{% 
	  \tt {\tiny [A$\rangle$GM]}  
	 ARGUMENT: {'Ich verstehe Ihren Wunsch, B66 zu behalten, aber C60 hat für mich einen signifikanten Wert (484) und sollte nicht einfach entfernt werden. Lassen Sie uns stattdessen über B10 sprechen - es hat nur 31 Aufwand aber 215 Wert für mich. Vielleicht können wir B10 zu unserem Paket hinzufügen, ohne andere wertvolle Gegenstände zu opfern?'}\\ \tt \\ \tt VORSCHLAG: {'C41', 'C15', 'C28', 'C48', 'B78', 'A35', 'A56', 'B66', 'C60', 'B49', 'A93', 'B10'}\\ \tt \\ \tt ARGUMENT: {'Mit diesem Vorschlag (Aufwand: 4108) überschreiten wir leider das Budget leicht. Wir müssen einen Gegenstand entfernen, um unter 4100 zu bleiben. Vielleicht könnten wir auf A35 verzichten? Oder haben Sie einen anderen Vorschlag, welcher Gegenstand entfernt werden könnte?'} 
	  } 
	   } 
	   } 
	 & & \\ 
 

    \theutterance \stepcounter{utterance}  

    & & \multicolumn{4}{p{0.6\linewidth}}{\cellcolor[rgb]{0.9,0.9,0.9}{%
	\makecell[{{p{\linewidth}}}]{% 
	  \tt {\tiny [GM$\rangle$B]}  
	 ARGUMENT: {'Ich verstehe Ihren Wunsch, B66 zu behalten, aber C60 hat für mich einen signifikanten Wert (484) und sollte nicht einfach entfernt werden. Lassen Sie uns stattdessen über B10 sprechen - es hat nur 31 Aufwand aber 215 Wert für mich. Vielleicht können wir B10 zu unserem Paket hinzufügen, ohne andere wertvolle Gegenstände zu opfern?'}\\ \tt \\ \tt VORSCHLAG: {'C41', 'C15', 'C28', 'C48', 'B78', 'A35', 'A56', 'B66', 'C60', 'B49', 'A93', 'B10'}\\ \tt \\ \tt ARGUMENT: {'Mit diesem Vorschlag (Aufwand: 4108) überschreiten wir leider das Budget leicht. Wir müssen einen Gegenstand entfernen, um unter 4100 zu bleiben. Vielleicht könnten wir auf A35 verzichten? Oder haben Sie einen anderen Vorschlag, welcher Gegenstand entfernt werden könnte?'} 
	  } 
	   } 
	   } 
	 & \\ 
 

    \theutterance \stepcounter{utterance}  

    & & & \multicolumn{4}{p{0.6\linewidth}}{\cellcolor[rgb]{1,0.85,0.72}{%
	\makecell[{{p{\linewidth}}}]{% 
	  \tt {\tiny [GM$\langle$B]}  
	 ARGUMENT: {'Sie haben recht, dass 4108 über dem Budget liegt. Anstatt A35 zu entfernen, das für Sie wertvoll ist, schlage ich vor, B05 stattdessen hinzuzufügen. Es hat 737 Aufwand aber nur 508 Wert für mich - nicht ideal, aber besser als A35 zu verlieren. So bleiben alle wichtigen Gegenstände erhalten.'}\\ \tt \\ \tt VORSCHLAG: {'C41', 'C15', 'C28', 'C48', 'B78', 'A35', 'A56', 'B66', 'C60', 'B49', 'A93', 'B05'}\\ \tt \\ \tt ARGUMENT: {'Dieser Vorschlag hat 4100 Aufwand genau - perfekt ausgeschöpftes Budget! Er enthält alle Ihre Prioritäten einschließlich A35 und B10-Ersatz B05, plus meine wertvollen C28, C48, A56, A93 und B66. Ein fairer Kompromiss, der beide Seiten berücksichtigt.'} 
	  } 
	   } 
	   } 
	  \\ 
 

    \theutterance \stepcounter{utterance}  

    & & \multicolumn{4}{p{0.6\linewidth}}{\cellcolor[rgb]{0.9,0.9,0.9}{%
	\makecell[{{p{\linewidth}}}]{% 
	  \tt {\tiny [A$\langle$GM]}  
	 ARGUMENT: {'Sie haben recht, dass 4108 über dem Budget liegt. Anstatt A35 zu entfernen, das für Sie wertvoll ist, schlage ich vor, B05 stattdessen hinzuzufügen. Es hat 737 Aufwand aber nur 508 Wert für mich - nicht ideal, aber besser als A35 zu verlieren. So bleiben alle wichtigen Gegenstände erhalten.'}\\ \tt \\ \tt VORSCHLAG: {'C41', 'C15', 'C28', 'C48', 'B78', 'A35', 'A56', 'B66', 'C60', 'B49', 'A93', 'B05'}\\ \tt \\ \tt ARGUMENT: {'Dieser Vorschlag hat 4100 Aufwand genau - perfekt ausgeschöpftes Budget! Er enthält alle Ihre Prioritäten einschließlich A35 und B10-Ersatz B05, plus meine wertvollen C28, C48, A56, A93 und B66. Ein fairer Kompromiss, der beide Seiten berücksichtigt.'} 
	  } 
	   } 
	   } 
	 & \\ 
 

    \theutterance \stepcounter{utterance}  

    & \multicolumn{4}{p{0.6\linewidth}}{\cellcolor[rgb]{0.8,1,0.9}{%
	\makecell[{{p{\linewidth}}}]{% 
	  \tt {\tiny [A$\rangle$GM]}  
	 ARGUMENT: {'Während ich die Berücksichtigung meiner Prioritäten schätze, ist B05 mit nur 508 Wert für mich im Vergleich zu seinem hohen Aufwand von 737 nicht optimal. B10 wäre mit 215 Wert bei nur 31 Aufwand eine viel effizientere Wahl für mich. Könnten wir stattdessen B10 anstelle von B05 aufnehmen?'}\\ \tt \\ \tt VORSCHLAG: {'C41', 'C15', 'C28', 'C48', 'B78', 'A35', 'A56', 'B66', 'C60', 'B49', 'A93', 'B10'}\\ \tt \\ \tt ARGUMENT: {'Dieser Vorschlag hat 4096 Aufwand - unter dem Budget, aber mit besserer Wert-Aufwand-Effizienz für mich. Wir behalten alle wichtigen Gegenstände beider Seiten und fügen B10 hinzu, das für mich wertvoller pro Aufwand ist als B05.'} 
	  } 
	   } 
	   } 
	 & & \\ 
 

    \theutterance \stepcounter{utterance}  

    & & \multicolumn{4}{p{0.6\linewidth}}{\cellcolor[rgb]{0.9,0.9,0.9}{%
	\makecell[{{p{\linewidth}}}]{% 
	  \tt {\tiny [GM$\rangle$B]}  
	 ARGUMENT: {'Während ich die Berücksichtigung meiner Prioritäten schätze, ist B05 mit nur 508 Wert für mich im Vergleich zu seinem hohen Aufwand von 737 nicht optimal. B10 wäre mit 215 Wert bei nur 31 Aufwand eine viel effizientere Wahl für mich. Könnten wir stattdessen B10 anstelle von B05 aufnehmen?'}\\ \tt \\ \tt VORSCHLAG: {'C41', 'C15', 'C28', 'C48', 'B78', 'A35', 'A56', 'B66', 'C60', 'B49', 'A93', 'B10'}\\ \tt \\ \tt ARGUMENT: {'Dieser Vorschlag hat 4096 Aufwand - unter dem Budget, aber mit besserer Wert-Aufwand-Effizienz für mich. Wir behalten alle wichtigen Gegenstände beider Seiten und fügen B10 hinzu, das für mich wertvoller pro Aufwand ist als B05.'} 
	  } 
	   } 
	   } 
	 & \\ 
 

    \theutterance \stepcounter{utterance}  

    & & & \multicolumn{4}{p{0.6\linewidth}}{\cellcolor[rgb]{1,0.85,0.72}{%
	\makecell[{{p{\linewidth}}}]{% 
	  \tt {\tiny [GM$\langle$B]}  
	 ARGUMENT: {'Ich verstehe Ihre Präferenz für B10, aber der Verlust von B05 (508 Wert für mich) ist erheblich. Lassen Sie uns einen Mittelweg finden: Behalten wir B05 bei und fügen B10 hinzu, indem wir einen weniger wertvollen Gegenstand für mich entfernen. Vielleicht C41? Es hat nur 138 Wert für mich bei 108 Aufwand.'}\\ \tt \\ \tt VORSCHLAG: {'C15', 'C28', 'C48', 'B78', 'A35', 'A56', 'B66', 'C60', 'B49', 'A93', 'B05', 'B10'}\\ \tt \\ \tt ARGUMENT: {'Dieser Vorschlag (4099 Aufwand) behält B05 für Sie und fügt B10 hinzu, während C41 entfernt wird. Obwohl C41 für mich wertvoll war, ist dies ein fairer Kompromiss, um Ihre Effizienzanliegen zu adressieren und gleichzeitig Ihren hohen Wert für B05 zu bewahren.'} 
	  } 
	   } 
	   } 
	  \\ 
 

    \theutterance \stepcounter{utterance}  

    & & \multicolumn{4}{p{0.6\linewidth}}{\cellcolor[rgb]{0.9,0.9,0.9}{%
	\makecell[{{p{\linewidth}}}]{% 
	  \tt {\tiny [A$\langle$GM]}  
	 ARGUMENT: {'Ich verstehe Ihre Präferenz für B10, aber der Verlust von B05 (508 Wert für mich) ist erheblich. Lassen Sie uns einen Mittelweg finden: Behalten wir B05 bei und fügen B10 hinzu, indem wir einen weniger wertvollen Gegenstand für mich entfernen. Vielleicht C41? Es hat nur 138 Wert für mich bei 108 Aufwand.'}\\ \tt \\ \tt VORSCHLAG: {'C15', 'C28', 'C48', 'B78', 'A35', 'A56', 'B66', 'C60', 'B49', 'A93', 'B05', 'B10'}\\ \tt \\ \tt ARGUMENT: {'Dieser Vorschlag (4099 Aufwand) behält B05 für Sie und fügt B10 hinzu, während C41 entfernt wird. Obwohl C41 für mich wertvoll war, ist dies ein fairer Kompromiss, um Ihre Effizienzanliegen zu adressieren und gleichzeitig Ihren hohen Wert für B05 zu bewahren.'} 
	  } 
	   } 
	   } 
	 & \\ 
 

    \theutterance \stepcounter{utterance}  

    & \multicolumn{4}{p{0.6\linewidth}}{\cellcolor[rgb]{0.8,1,0.9}{%
	\makecell[{{p{\linewidth}}}]{% 
	  \tt {\tiny [A$\rangle$GM]}  
	 ARGUMENT: {'Ich sehe, dass Sie bereit sind, auf C41 zu verzichten, was ich zu schätzen weiß. Allerdings hat C41 mit 138 Wert bei nur 108 Aufwand eine bessere Effizienz für mich als B05. Könnten wir stattdessen beide behalten und einen anderen Gegenstand mit geringerer Effizienz entfernen? Vielleicht B66 (461 Wert bei 855 Aufwand) oder A35 (389 Wert bei 533 Aufwand)?'}\\ \tt \\ \tt VORSCHLAG: {'C41', 'C15', 'C28', 'C48', 'B78', 'A56', 'C60', 'B49', 'A93', 'B05', 'B10'}\\ \tt \\ \tt ARGUMENT: {'Dieser Vorschlag hat 4087 Aufwand und entfernt sowohl B66 als auch A35. Er behält alle hochwertigen Gegenstände für uns beide, einschließlich Ihres B05 und meines C41, plus B10 für zusätzliche Effizienz. Ein fairer Kompromiss, der die Gesamteffizienz des Pakets verbessert.'} 
	  } 
	   } 
	   } 
	 & & \\ 
 

    \theutterance \stepcounter{utterance}  

    & & \multicolumn{4}{p{0.6\linewidth}}{\cellcolor[rgb]{0.9,0.9,0.9}{%
	\makecell[{{p{\linewidth}}}]{% 
	  \tt {\tiny [GM$\rangle$B]}  
	 ARGUMENT: {'Ich sehe, dass Sie bereit sind, auf C41 zu verzichten, was ich zu schätzen weiß. Allerdings hat C41 mit 138 Wert bei nur 108 Aufwand eine bessere Effizienz für mich als B05. Könnten wir stattdessen beide behalten und einen anderen Gegenstand mit geringerer Effizienz entfernen? Vielleicht B66 (461 Wert bei 855 Aufwand) oder A35 (389 Wert bei 533 Aufwand)?'}\\ \tt \\ \tt VORSCHLAG: {'C41', 'C15', 'C28', 'C48', 'B78', 'A56', 'C60', 'B49', 'A93', 'B05', 'B10'}\\ \tt \\ \tt ARGUMENT: {'Dieser Vorschlag hat 4087 Aufwand und entfernt sowohl B66 als auch A35. Er behält alle hochwertigen Gegenstände für uns beide, einschließlich Ihres B05 und meines C41, plus B10 für zusätzliche Effizienz. Ein fairer Kompromiss, der die Gesamteffizienz des Pakets verbessert.'} 
	  } 
	   } 
	   } 
	 & \\ 
 

    \theutterance \stepcounter{utterance}  

    & & & \multicolumn{4}{p{0.6\linewidth}}{\cellcolor[rgb]{1,0.85,0.72}{%
	\makecell[{{p{\linewidth}}}]{% 
	  \tt {\tiny [GM$\langle$B]}  
	 ARGUMENT: {'Ich kann Ihre Effizienzargumente nachvollziehen, aber B66 (461 Wert) und A35 (389 Wert) sind zu wertvoll für mich, um sie beide zu opfern. Lassen Sie uns stattdessen nur A35 entfernen, da es einen geringeren Wert pro Aufwand hat als B66, und beide B05 und C41 behalten.'}\\ \tt \\ \tt VORSCHLAG: {'C41', 'C15', 'C28', 'C48', 'B78', 'A56', 'B66', 'C60', 'B49', 'A93', 'B05', 'B10'}\\ \tt \\ \tt ARGUMENT: {'Mit 4091 Aufwand bleibt dieser Vorschlag unter dem Budget und bietet die beste Balance: Ihre effizienten Items B05 und B10, mein wertvolles C41, und wir behalten B66, das für mich wichtiger ist als A35. Ein fairer Kompromiss, der beide Seiten zufriedenstellt.'} 
	  } 
	   } 
	   } 
	  \\ 
 

    \theutterance \stepcounter{utterance}  

    & & \multicolumn{4}{p{0.6\linewidth}}{\cellcolor[rgb]{0.9,0.9,0.9}{%
	\makecell[{{p{\linewidth}}}]{% 
	  \tt {\tiny [A$\langle$GM]}  
	 ARGUMENT: {'Ich kann Ihre Effizienzargumente nachvollziehen, aber B66 (461 Wert) und A35 (389 Wert) sind zu wertvoll für mich, um sie beide zu opfern. Lassen Sie uns stattdessen nur A35 entfernen, da es einen geringeren Wert pro Aufwand hat als B66, und beide B05 und C41 behalten.'}\\ \tt \\ \tt VORSCHLAG: {'C41', 'C15', 'C28', 'C48', 'B78', 'A56', 'B66', 'C60', 'B49', 'A93', 'B05', 'B10'}\\ \tt \\ \tt ARGUMENT: {'Mit 4091 Aufwand bleibt dieser Vorschlag unter dem Budget und bietet die beste Balance: Ihre effizienten Items B05 und B10, mein wertvolles C41, und wir behalten B66, das für mich wichtiger ist als A35. Ein fairer Kompromiss, der beide Seiten zufriedenstellt.'} 
	  } 
	   } 
	   } 
	 & \\ 
 

    \theutterance \stepcounter{utterance}  

    & \multicolumn{4}{p{0.6\linewidth}}{\cellcolor[rgb]{0.8,1,0.9}{%
	\makecell[{{p{\linewidth}}}]{% 
	  \tt {\tiny [A$\rangle$GM]}  
	 ZUSTIMMUNG: {'C41', 'C15', 'C28', 'C48', 'B78', 'A56', 'B66', 'C60', 'B49', 'A93', 'B05', 'B10'}\\ \tt \\ \tt ARGUMENT: {'Ich akzeptiere diesen Vorschlag, da er eine gute Balance zwischen unseren Interessen findet und unter dem Budget bleibt. Der Vorschlag enthält meine wichtigsten Gegenstände wie C28, C48, A56 und A93, während auch Ihre Prioritäten B05 und B10 berücksichtigt werden.'} 
	  } 
	   } 
	   } 
	 & & \\ 
 

    \theutterance \stepcounter{utterance}  

    & & & \multicolumn{2}{p{0.3\linewidth}}{\cellcolor[rgb]{0.95,0.95,0.95}{%
	\makecell[{{p{\linewidth}}}]{% 
	  \tt {\tiny [GM$|$GM]}  
	 game successful 
	  } 
	   } 
	   } 
	 & & \\ 
 

    \theutterance \stepcounter{utterance}  

    & & & \multicolumn{2}{p{0.3\linewidth}}{\cellcolor[rgb]{0.95,0.95,0.95}{%
	\makecell[{{p{\linewidth}}}]{% 
	  \tt {\tiny [GM$|$GM]}  
	 end game 
	  } 
	   } 
	   } 
	 & & \\ 
 

\end{supertabular}
}

\end{document}
